\section{The result and the plan}\label{sec:intro}

\subsection{The result}

Throughout, $\log$ is the binary logarithm, and $\lgk{j}$ is its $j$-fold iterate. For a graph
$G$, an \emph{even subgraph} is a set of edges that meets every vertex an even number of times.
With symmetric difference as addition, the even subgraphs form a vector space $\Z(G)$ over
$\mathbb{F}_2$. Its dimension is
\[
  \beta(G)=|E(G)|-|V(G)|+\comp(G),
\]
where $\comp(G)$ is the number of components. We call $\beta(G)$ the \emph{rank} of $G$; other
names are cyclomatic number and cycle rank. A \emph{cycle} is a connected $2$-regular subgraph. We
identify a subgraph without isolated vertices with its edge set. Then a cycle is a nonzero element
of $\Z(G)$, so a graph of rank $r$ has at most $2^{r}-1$ cycles.

Let $\psi(r)$ be the largest number of cycles in a graph of rank $r$.
\begin{itemize}[leftmargin=1.4em]
\item Entringer and Slater~\cite{ES} proved that $2^{r-1}<\psi(r)<2^{r}$ for $r\ge3$, and they
conjectured that $\psi(r)\sim2^{r-1}$.
\item Aldred and Thomassen~\cite{AT} proved that $\psi(r)\le\frac{15}{16}2^{r}$.
\item Alahmadi, Aldred, dela Cruz, Sol\'e and Thomassen~\cite{AACST} proved that a connected
graph with $p$ vertices and more than $2p+3\log(3p)$ edges has at most $2^{r-1}$ cycles.
\item Horak and Kim~\cite{HK} found cubic graphs of rank $r$ with $H(r)$ cycles, where
\begin{equation}\label{eq:H}
  H(2q+1)=2^{2q}+\bigl(2q+\tfrac12\bigr)2^{q}-3q,\qquad
  H(2q+2)=2^{2q+1}+\bigl(2q+\tfrac92\bigr)2^{q}-3q-4 .
\end{equation}
\end{itemize}
For odd rank $r=2q+1$ the graph of Horak and Kim is a ring of $q$ copies of $K_{2,2}$, each joined
to the next by two edges. For even rank, one more edge is inserted across the ring
(Figure~\ref{fig:ring}). We call this graph $R_r$.

\begin{figure}[ht]
\centering
\begin{tikzpicture}
  \def\ri{1.15}\def\ro{1.95}\def\dl{15}
  \foreach \i/\th in {1/18,2/-54,3/-126,4/162}{
    \coordinate (li\i) at (\th+\dl:\ri);
    \coordinate (lo\i) at (\th+\dl:\ro);
    \coordinate (ri\i) at (\th-\dl:\ri);
    \coordinate (ro\i) at (\th-\dl:\ro);
    \draw[cool] (li\i)--(ri\i) (lo\i)--(ro\i) (li\i)--(ro\i) (lo\i)--(ri\i);
  }
  \coordinate (xi) at (90:\ri);
  \coordinate (xo) at (90:\ro);
  \draw[hot] (xi)--(xo);
  \draw[edge] (ri4)--(xi) (ro4)--(xo) (xi)--(li1) (xo)--(lo1);
  \draw[edge] (ri1)--(li2) (ro1)--(lo2) (ri2)--(li3) (ro2)--(lo3) (ri3)--(li4) (ro3)--(lo4);
  \foreach \i in {1,2,3,4}{\foreach \n in {li,lo,ri,ro}{\node[vtx] at (\n\i) {};}}
  \node[vtx] at (xi) {};
  \node[vtx] at (xo) {};
  \node[note, text=warm, anchor=west] at (2.75,1.75) {the further edge (even rank only)};
  \draw[warm, line width=0.4pt] (2.7,1.75) -- (0.12,1.6);
  \node[note, text=accent, anchor=west] at (2.75,-0.55) {a copy of $K_{2,2}$};
  \draw[accent, line width=0.4pt] (2.7,-0.55) -- (1.7,-0.62);
  \node[note, anchor=west] at (2.75,0.6) {two edges join consecutive copies};
  \draw[slate, line width=0.4pt] (2.7,0.6) -- (1.62,0.98);
\end{tikzpicture}
\caption{The graph $R_{10}$: a ring of four copies of $K_{2,2}$ and one further edge. It has rank
$10$ and $H(10)=696$ cycles.}
\label{fig:ring}
\end{figure}

\begin{maintheorem}{the maximum number of cycles}\label{thm:exact}
There is an integer $R$ such that $\psi(r)=H(r)$ for every $r\ge R$.
\end{maintheorem}

In particular $\psi(r)=2^{r-1}+(1+o(1))\,r\,2^{\lfloor(r-1)/2\rfloor}$, so $\psi(r)\sim2^{r-1}$.
The proof gives no value for $R$.

\begin{fineprint}
The theorem is false for small $r$. For example, take the vertices $u_i,v_i$ with
$i\in\mathbb{Z}/26$ and the edges $u_iu_{i+1}$, $u_iv_i$ and $v_iv_{i+10}$. This generalised
Petersen graph has rank $27$, and a direct count by computer gives $67\,959\,973$ cycles, while
$H(27)=67\,325\,913$.
\end{fineprint}

\subsection{The three cases of the proof}

The number of cycles does not change when we delete a vertex of degree at most one, or replace a
vertex of degree two and its two edges by one edge. A vertex of degree at least four can be split
into vertices of degree three, and no cycle is lost. So it suffices to count the cycles of a
connected multigraph $G$ of rank $r$ in which all vertices have degree three
(Lemma~\ref{lem:cubic}). The graphs $R_r$ show that $\psi(r)\ge H(r)$.

For the upper bound there are three cases (Figure~\ref{fig:plan}). A \emph{cyclic cut} is a set of
edges whose deletion leaves two parts that each contain a cycle. The number $H(r)-2^{r-1}$ is of
order $r\,2^{r/2}$, so the three bounds give the theorem (Section~\ref{sec:proof}).

\begin{figure}[ht]
\centering
\begin{tikzpicture}[
  ask/.style={draw=accent, line width=0.7pt, rounded corners=2pt, fill=white,
    font=\sffamily\footnotesize, text width=4.6cm, align=center, inner sep=4.5pt},
  res/.style={fill=accent!7, rounded corners=2pt, font=\sffamily\footnotesize,
    text width=6.3cm, align=left, inner sep=4.5pt},
  ar/.style={-{Stealth[length=4.5pt]}, line width=0.6pt, draw=slate},
  lab/.style={font=\sffamily\scriptsize, text=slate, inner sep=1.5pt},
  node distance=5.2mm]
  \node[ask] (q2) {Does $G$ have a cyclic cut of at most three edges?};
  \node[ask, below=of q2] (q3) {Does $G$ have a cycle of length four?};
  \node[res, right=13mm of q2] (o2) {at most $2^{r-1}+O(2^{r/2})$ cycles\\ \textcolor{slate}{Proposition~\ref{prop:threecut}}};
  \node[res, right=13mm of q3] (o3) {a ring, with at most $H(r)$ cycles, or fewer than $2^{r-1}$ cycles\\ \textcolor{slate}{Theorem~\ref{thm:fourcut}, Corollary~\ref{cor:ringmax}}};
  \node[res, below=of q3, text width=4.6cm, align=center] (o4) {girth at least five, and fewer than $2^{r-1}$ cycles\\ \textcolor{slate}{Theorem~\ref{thm:half}}};
  \foreach \a/\b in {q2/o2,q3/o3} \draw[ar] (\a) -- node[lab, above] {yes} (\b);
  \foreach \a/\b in {q2/q3,q3/o4} \draw[ar] (\a) -- node[lab, right] {no} (\b);
\end{tikzpicture}
\caption{The three cases for a connected cubic multigraph $G$ of large rank $r$.}
\label{fig:plan}
\end{figure}

\subsection{Outside parts and the forest estimate}

Fix a set $P$ of vertices of $G$. For an even subgraph $\xi$ we call the set of its edges without
an end in $P$ its \emph{outside part} (Figure~\ref{fig:outside}). If a cycle meets $P$, then its
outside part is a proper subset of a cycle, so it is a forest.

\begin{figure}[ht]
\centering
\begin{tikzpicture}[x=1cm,y=1cm]
  % the set P: a hexagon
  \foreach \i in {0,...,5} \coordinate (w\i) at ({60*\i+30}:0.9);
  \draw[cool] (w0)--(w1)--(w2)--(w3)--(w4)--(w5)--cycle;
  % outside vertices
  \coordinate (a) at (2.1,1.05);
  \coordinate (b) at (3.3,0.25);
  \coordinate (c) at (2.1,-1.05);
  \coordinate (d) at (-2.2,0.9);
  \coordinate (e) at (-3.0,-0.1);
  \coordinate (f) at (-2.0,-0.9);
  \draw[edge] (w3)--(d) (w2)--(d) (w3)--(f);
  \draw[edge] (d)--(e)--(f) (a)--(c);
  % the cycle C: w0 - a - b - c - w5 - w0 ; outside part thick
  \draw[hot] (a)--(b)--(c);
  \draw[hot, dash pattern=on 2.2pt off 1.6pt] (w5)--(w0) (w0)--(a) (c)--(w5);
  \foreach \p in {w0,w1,w2,w3,w4,w5,a,b,c,d,e,f} \node[vtx] at (\p) {};
  \node[note, text=accent] at (0,0) {$P$};
  \node[note, text=warm, anchor=west] at (3.45,0.25) {the outside part of the cycle $C$: a path};
\end{tikzpicture}
\caption{A set $P$ of six vertices and a cycle $C$ that meets $P$. The edges of $C$ with an end
in $P$ are dashed. The other edges of $C$ form its outside part, which is a path. An even
subgraph whose outside part contains a cycle is not a cycle that meets $P$.}
\label{fig:outside}
\end{figure}

\begin{idea}{cycles and outside parts}
Let $P$ be a small set of vertices such that $G[P]$ is connected. By the \emph{forest estimate}
of the companion paper~\cite{Forest} (Theorem~\ref{thm:forest}), the outside part is a forest for
at most $\half+o(1)$ of all even subgraphs. So at most $(\half+o(1))2^{r}$ cycles meet $P$. A cycle that avoids $P$ is a cycle of
$G-P$, and if $G$ has no bridge, the rank of $G-P$ is at most $r-1-|P|/2$
(Lemma~\ref{lem:marked}). This gives $\psi(r)\le(\half+o(1))2^{r}$ (Corollary~\ref{cor:cycles}).
\end{idea}

In Cases~1 and~2 a cycle consists of paths in the parts of $G$. By the forest estimate,
a part with more than half of the possible paths between two of its ends, by a fixed margin, is a
ladder up to a bounded error, and it has few cycles (Section~\ref{sec:poles}). In Case~2 we
also put a small tree, a \emph{cap}, in the place of the cycle of length four, and apply the
bound $(\half+o(1))2^{r}$ to the graph that arises (Lemma~\ref{lem:cap}). In Case~3 the
forest estimate is not sufficient, because we need a bound below $\half$ and its error term is
positive. In this case a second-moment argument gives two disjoint pieces of $G$ with small
boundaries, and each piece contributes a factor of at most $\frac35$ to the proportion of
forests (Section~\ref{sec:half}).

\subsection{Authorship and verification}\label{sec:status}

KNT is a person. Astra and Fable~5.2 are the language models GPT-6 Astra and Claude Fable~5.2.
KNT directed the work at each step. KNT obtained the forest estimate in September 2026 with
Astra, for a problem of Erd\H{o}s on pancyclic graphs~\cite{KNT}; its proof is
in~\cite{Forest}. For the number of cycles, Astra found the path structure of
Section~\ref{sec:poles}, the gap theorem, the count for rings, the cap, and the first proof of
the last case. Fable~5.2 found the lemma on layers, the chain of separators and the pieces at
girth five, and wrote this text. The present forms of the proofs contain steps of both models.

Theorem~\ref{thm:exact} has a formal proof in Lean~4. The statement
\lean{Erdos1016.ExactMaximum.exactMaximum_threeCases} says that there is a number $R$ such that
\lean{maxCycles r = horakKim r} for every $r\ge R$. Here \lean{maxCycles r} is the largest number
of cycles of a simple graph with cyclomatic number $r$, and \lean{horakKim r} is $H(r)$. The Lean
kernel accepts the proof with the axioms \lean{propext}, \lean{Classical.choice} and
\lean{Quot.sound} only. The files are in~\cite{Repo}, at commit \texttt{b7f9fc9}.

The formal proof has the three cases of Section~\ref{sec:proof}, and it proves
Theorem~\ref{thm:half} at girth five. Some of its steps are not those of the text: a cut of
size four has two arbitrary sides, Lemma~\ref{lem:fullfive} has a proof without ladders,
and the end of the proof of Theorem~\ref{thm:half} uses short cycles at many scales. So the
kernel check
covers the theorem, and it does not cover the text of this paper. No referee has read the text.
